% Diagramma dei casi d'uso - HomeStock
% Compila con pdfLaTeX (funziona anche su Overleaf)
\documentclass[tikz,border=12pt]{standalone}
\usepackage[utf8]{inputenc}
\usepackage[T1]{fontenc}
\usepackage[italian]{babel}
\usepackage{lmodern}
\usetikzlibrary{shapes.geometric, arrows.meta, backgrounds}

\definecolor{hsteal}{HTML}{0E8F89}
\definecolor{hsink}{HTML}{16211E}
\definecolor{hsgrey}{HTML}{66706B}

\tikzset{
  % caso d'uso
  uc/.style={ellipse, draw=hsteal, line width=1pt, fill=white,
             minimum width=4.6cm, minimum height=1.1cm,
             text width=3.7cm, align=center, font=\small},
  % attore: nodo invisibile a cui si collegano le linee
  attore/.style={inner sep=0pt, minimum width=0.8cm, minimum height=1.5cm},
  % associazione attore - caso d'uso
  assoc/.style={draw=hsgrey, line width=0.6pt},
  % include / extend
  rel/.style={draw=hsink, line width=0.9pt, dash pattern=on 7pt off 4pt,
              -{Straight Barb[length=3.5mm, width=3mm]}},
  etichetta/.style={font=\footnotesize\itshape, fill=white, inner sep=1.5pt},
  % omino stilizzato
  omino/.pic={
    \draw[hsink, line width=1pt]
      (0,0.52) circle (0.2)
      (0,0.32) -- (0,-0.2)
      (-0.36,0.14) -- (0.36,0.14)
      (0,-0.2) -- (-0.28,-0.68)
      (0,-0.2) -- (0.28,-0.68);
  }
}

\newcommand{\attore}[3]{% nome nodo, posizione, etichetta
  \node[attore] (#1) at (#2) {};
  \pic at (#1) {omino};
  \node[font=\small\bfseries, below=0.78cm, align=center] at (#1) {#3};
}

\begin{document}
\begin{tikzpicture}

% ---------- confine del sistema ----------
\draw[hsink, line width=1pt, rounded corners=4pt] (-3.0,1.6) rectangle (10.2,-15.2);
\node[font=\large\bfseries] at (3.6,1.05) {HomeStock};

% ---------- casi d'uso: colonna principale ----------
\node[uc] (reg)   at (0,   0.00) {Registrazione e accesso};
\node[uc] (loc)   at (0,  -1.55) {Organizzazione per luogo};
\node[uc] (man)   at (0,  -3.10) {Aggiunta manuale prodotti};
\node[uc] (exp)   at (0,  -4.65) {Date di scadenza};
\node[uc] (inv)   at (0,  -6.20) {Visualizzazione inventario};
\node[uc] (qta)   at (0,  -7.75) {Gestione quantità};
\node[uc] (spesa) at (0,  -9.30) {Lista della spesa};
\node[uc] (ric)   at (0, -10.85) {Suggerimento ricette};
\node[uc] (cond)  at (0, -12.40) {Condivisione familiare};
\node[uc] (prem)  at (0, -13.95) {Premium};

% ---------- casi d'uso: estensioni e inclusioni ----------
\node[uc] (scan)  at (7.2,  -3.10) {Scansione codice a barre};
\node[uc] (notif) at (7.2,  -4.65) {Avvisi e notifiche};
\node[uc] (cerca) at (7.2,  -6.20) {Ricerca e ordinamento};
\node[uc] (trasf) at (7.2,  -9.30) {Trasferimento in dispensa};
\node[uc] (ingr)  at (7.2, -10.85) {Ingredienti disponibili e mancanti};
\node[uc] (prio)  at (7.2, -12.40) {Priorità ai prodotti in scadenza};
\node[uc] (stat)  at (7.2, -13.95) {Statistiche e spreco};

% ---------- attori ----------
\attore{visitatore}{-6.2, -0.4}{Visitatore}
\attore{utente}{-6.2, -5.6}{Utente}
\attore{membro}{-6.2, -9.6}{Membro\\famiglia}
\attore{premium}{12.4, -13.95}{Utente\\Premium}

% ---------- associazioni (sotto le ellissi) ----------
\begin{pgfonlayer}{background}
\draw[assoc] (visitatore) -- (reg);

\foreach \u in {reg, loc, man, exp, inv, qta, spesa, ric, cond, prem}
  \draw[assoc] (utente) -- (\u);

\foreach \u in {man, inv, qta, spesa, cond}
  \draw[assoc] (membro) -- (\u);

\draw[assoc] (premium) -- (stat);
\end{pgfonlayer}

% ---------- extend: dall'estensione verso il caso base ----------
\draw[rel] (scan)  -- node[etichetta, above] {«extend»} (man);
\draw[rel] (notif) -- node[etichetta, above] {«extend»} (exp);
\draw[rel] (cerca) -- node[etichetta, above] {«extend»} (inv);
\draw[rel] (trasf) -- node[etichetta, above] {«extend»} (spesa);

% ---------- include: dal caso base verso quello incluso ----------
\draw[rel] (ric) -- node[etichetta, above] {«include»} (ingr);
\draw[rel] (ric) -- node[etichetta, sloped, above] {«include»} (prio);

\end{tikzpicture}
\end{document}
